% !TEX root = ../../../main.tex
\section{Intermediate value theorem}
\label{sec:analysis-i:intermediate-value-theorem}

% BEGIN TUFTE PLACEHOLDER: supplied sample, lines 765--783.
\begingroup\sloppy
% Replace this entire specimen block with reviewed course notes.
If, after using the \doccmd{morefloats} command twice, you continue to get the
\texttt{Too many unprocessed floats} error, there are a couple things you can
do.

The \doccmddef{FloatBarrier} command will immediately process all the floats
before typesetting more material.  Since \doccmd{FloatBarrier} will start a new
paragraph, you should place this command at the beginning or end of a
paragraph.

The \doccmddef{clearpage} command will also process the floats before
continuing, but instead of starting a new paragraph, it will start a new page.

You can also try moving your floats around a bit: move a figure or table to the
next page or reduce the number of sidenotes.  (Each sidenote actually uses
\emph{two} slots.)

After the floats have placed, \LaTeX{} will mark those slots as unused so they
are available for the next page to be composed.

\lipsum[7]
\par\endgroup
% END TUFTE PLACEHOLDER
